\documentclass[10pt]{beamer}
\usetheme{Madrid}
\usepackage[T1]{fontenc}
\usepackage[utf8]{inputenc}
\usepackage[brazil]{babel}
\usepackage{lmodern}
\usepackage{hyperref}
\usepackage{listings}
\usepackage{xcolor}
\setbeamertemplate{navigation symbols}{}

\lstset{
  language=C++,
  basicstyle=\ttfamily\small,
  keywordstyle=\color{blue},
  commentstyle=\color{gray},
  stringstyle=\color{red},
  breaklines=true,
  showstringspaces=false,
}

\title{Máquinas de Estado e Flags}
\subtitle{Curso de Microcontroladores}
\author{}
\date{}

\begin{document}

\begin{frame}
  \titlepage
\end{frame}

\begin{frame}{O problema com `delay()`}
  \begin{itemize}
    \item `delay()` é bastante utilizado quando precisamos esperar para fazer uma ação.
    \item Porém, `delay()` bloqueia o `loop()` e paralisa o programa.
    \item Com o programa paralisado, você perde o controle sobre variáveis, funções e etc.
    \item Uma máquina de estado pode deixar o `loop()` sempre rodando, mas também mudar uma ação após um determinado tempo.
    \item É essencial no seguidor, pois os sensores sempre precisam estar funcionando para a tomada de decisão.
  \end{itemize}
\end{frame}

\begin{frame}[fragile]{Definindo estados com `enum`}
\begin{lstlisting}
enum Estado {
  PARADO,
  ANDANDO
};

Estado estadoAtual = PARADO;
bool mudou = false;
\end{lstlisting}

  \begin{itemize}
    \item `enum` cria um tipo com nomes simbólicos.
    \item Muito mais legível que usar números (0, 1, 2...).
    \item `mudou` é a flag que controla transições.
  \end{itemize}
\end{frame}

\begin{frame}[fragile]{Como usar no `loop()`}
\begin{lstlisting}
void loop() {
  if (!mudou) {
    if (estadoAtual == PARADO) {
      estadoAtual = ANDANDO;
      digitalWrite(LED_PIN, HIGH);
      mudou = true;
    }
  }
  delay(1000);
  mudou = false;
}
\end{lstlisting}
\end{frame}

\begin{frame}{Entendendo a lógica}
  \begin{itemize}
    \item `if (!mudou)` verifica se já mudou de estado.
    \item Dentro, você executa a ação correspondente ao estado.
    \item `mudou = true` marca que a transição aconteceu.
    \item No fim do `loop()`, `mudou` volta a `false` para a próxima oportunidade.
  \end{itemize}
\end{frame}

\begin{frame}{Vantagem para robôs}
  \begin{itemize}
    \item O `loop()` roda sempre, sensores são lidos continuamente.
    \item Estados definem claramente o que o robô está fazendo.
    \item Fácil adicionar novas transições e comportamentos.
  \end{itemize}
\end{frame}

\begin{frame}{Próximos passos}
  \begin{itemize}
    \item Teste o código com um LED que liga/desliga.
    \item Adicione um terceiro estado: `RETORNANDO`.
    \item Use um botão para controlar as transições de estado.
    \item Estude `millis()` para temporizações sem bloquear.
  \end{itemize}
\end{frame}

\end{document}
